\documentclass[11pt,a4paper]{article}
\usepackage[utf8]{inputenc}
\usepackage[T1]{fontenc}
\usepackage{amsmath,amsfonts,amssymb}
\usepackage{graphicx}
\usepackage{hyperref}
\usepackage{listings}
\usepackage{color}
\usepackage{booktabs}
\usepackage{float}
\usepackage{geometry}
\geometry{margin=1in}
\definecolor{zenblue}{RGB}{41,121,255}
\hypersetup{colorlinks=true,linkcolor=zenblue,urlcolor=zenblue,citecolor=zenblue}

\title{\textbf{Knowledge Distillation for Efficient Zen Models:\\
Semantic-Preserving Multi-Teacher Distillation}\\
\large Technical Report v2025.07}
\author{Antje Worring, Zach Kelling \\ Zen LM Research Team\\
\texttt{research@zenlm.org}}
\date{July 2025}

\begin{document}
\maketitle

\begin{abstract}
Knowledge distillation compresses large ``teacher'' models into smaller, faster
``student'' models while retaining as much task performance as possible. We present
a comprehensive distillation framework for the Qwen3-based Zen models (dense
0.6B/4B/8B/32B and the Qwen3-30B-A3B mixture-of-experts variant, all Apache-2.0),
covering offline distillation, online distillation, progressive distillation, and
task-specific distillation. Our central innovation is \emph{semantic-preserving
distillation (SPD)}, which augments the standard KL-divergence objective with a semantic
alignment loss that matches student and teacher representations in a shared anchor
embedding space. We further introduce \emph{multi-teacher distillation}, where multiple
teacher models of different specializations contribute to a single student. SPD closes a
substantially larger fraction of the teacher--student performance gap than standard
KL-divergence distillation, while producing students that are several times faster at
inference. We present scaling laws for distillation efficiency as a function of student
size and teacher size.
\end{abstract}

\tableofcontents
\newpage

%% -----------------------------------------------------------------------
\section{Introduction}
\label{sec:intro}
%% -----------------------------------------------------------------------

Deploying 32B-parameter models at scale demands significant compute: a single GPU serves
far fewer tokens/second for a 32B model in FP16 than for a 4B or 8B model. For
cost-sensitive applications (consumer products, edge devices, batch processing
pipelines), a well-distilled 4B or 8B student can provide near-teacher performance at a
fraction of the inference cost.

Knowledge distillation \cite{hinton2015distilling} trains a student model $S_\theta$
to mimic a teacher model $T$ by minimizing the KL divergence between their output
distributions:
\begin{equation}
    \mathcal{L}_{\text{KD}} = \text{KL}(p_T(\cdot \mid x) \,\|\, p_S(\cdot \mid x))
    \label{eq:kd}
\end{equation}

Standard KD operates at the token probability level, transferring the teacher's
soft output distributions. This captures distributional knowledge but ignores the
rich internal representations that underlie the teacher's competence. SPD augments
KD with an intermediate representation alignment loss in the semantic anchor space,
closing an additional portion of the teacher-student gap beyond what distributional
KD alone achieves.

\subsection{Distillation Taxonomy}

We distinguish four distillation regimes:

\begin{itemize}
    \item \textbf{Offline distillation}: teacher outputs precomputed and stored; student
          trains on static soft labels. Highest throughput, lowest cost.
    \item \textbf{Online distillation}: teacher runs synchronously during student
          training, providing dynamic soft labels. Higher quality, higher compute cost.
    \item \textbf{Progressive distillation}: iterative chain $T \to S_1 \to S_2 \to \ldots$,
          each stage compressing by a fixed factor. Enables extreme compression ratios.
    \item \textbf{Task-specific distillation}: distillation on domain-specific data with
          task-adapted teacher soft labels. Maximizes student performance on target tasks.
\end{itemize}

%% -----------------------------------------------------------------------
\section{Semantic-Preserving Distillation}
\label{sec:spd}
%% -----------------------------------------------------------------------

\subsection{Motivation}

Standard KD transfers distributional knowledge from teacher to student but does not
explicitly align their internal representations. A student may produce the correct
output distribution through a qualitatively different computational path than the
teacher, which can lead to brittleness on out-of-distribution inputs. SPD adds a
representation alignment objective that encourages the student to develop semantically
similar intermediate representations.

\subsection{SPD Loss}

Let $h_T^\ell(x) \in \mathbb{R}^{d_T}$ and $h_S^\ell(x) \in \mathbb{R}^{d_S}$ denote
the hidden states of teacher and student at layer $\ell$ (matched by relative depth).
SPD projects both to the shared anchor embedding space:
\begin{equation}
    \tilde{h}_T^\ell = \mathbf{P}_T h_T^\ell(x), \quad
    \tilde{h}_S^\ell = \mathbf{P}_S h_S^\ell(x)
    \label{eq:projection}
\end{equation}
where $\mathbf{P}_T \in \mathbb{R}^{d_A \times d_T}$ and
$\mathbf{P}_S \in \mathbb{R}^{d_A \times d_S}$ are alignment projection matrices
learned jointly with the student. The SPD alignment loss is:
\begin{equation}
    \mathcal{L}_{\text{align}} = \frac{1}{L}\sum_{\ell=1}^{L}
    \bigl(1 - \cos(\tilde{h}_T^\ell, \tilde{h}_S^\ell)\bigr)
    \label{eq:align_loss}
\end{equation}

The full SPD training objective is:
\begin{equation}
    \mathcal{L}_{\text{SPD}} = (1-\alpha)\mathcal{L}_{\text{CE}}(y, p_S)
    + \alpha \mathcal{L}_{\text{KD}}(p_T, p_S)
    + \beta \mathcal{L}_{\text{align}}
    \label{eq:spd_total}
\end{equation}
where $\alpha$ controls the blend of ground truth and teacher labels, and $\beta$
controls semantic alignment strength. Default: $\alpha = 0.5$, $\beta = 0.1$.

\subsection{Layer Matching Strategy}

Teacher and student have different depths ($L_T > L_S$). We match layers by relative
position: student layer $\ell$ corresponds to teacher layer $\lfloor \ell \cdot L_T / L_S \rfloor$.
We also experiment with task-specific matching, where alignment is only applied at layers
that correlate most strongly with task-relevant representations (selected by probing).

%% -----------------------------------------------------------------------
\section{Multi-Teacher Distillation}
\label{sec:multi_teacher}
%% -----------------------------------------------------------------------

\subsection{Motivation}

The Zen family includes specialized variants: a code-focused model, a mathematics-tuned
model, and a general-purpose model (each a fine-tune of a Qwen3-based Zen checkpoint).
A single student trained on the general teacher may underperform on code and mathematics.
Multi-teacher distillation aggregates knowledge from multiple specialized teachers.

\subsection{Multi-Teacher Loss}

Given $K$ teacher models $\{T_k\}_{k=1}^K$, the multi-teacher distillation loss is:
\begin{equation}
    \mathcal{L}_{\text{MT}} = (1-\alpha)\mathcal{L}_{\text{CE}} +
    \alpha \sum_{k=1}^{K} w_k(x) \cdot \mathcal{L}_{\text{KD}}(p_{T_k}, p_S) +
    \beta \sum_{k=1}^{K} w_k(x) \cdot \mathcal{L}_{\text{align}}^{(k)}
    \label{eq:mt_loss}
\end{equation}

where $w_k(x) \in [0,1]$ is a domain relevance weight for teacher $k$ on input $x$:
\begin{equation}
    w_k(x) = \text{softmax}\!\left(\frac{z_k(x)}{\tau}\right), \quad
    z_k(x) = \cos(e(x), c_k)
    \label{eq:teacher_weight}
\end{equation}
with $e(x)$ the input embedding and $c_k$ the centroid of teacher $k$'s training domain
in embedding space. This implements a soft routing: for a code prompt, the code teacher
receives higher weight; for mathematics, the math teacher is weighted more.

%% -----------------------------------------------------------------------
\section{Distillation Variants}
\label{sec:variants}
%% -----------------------------------------------------------------------

\subsection{Offline Distillation}

In offline distillation, teacher outputs $\{p_T(y \mid x)\}$ are precomputed for the
entire training corpus and stored as soft labels. The student trains on these stored
distributions using $\mathcal{L}_{\text{SPD}}$ with the alignment term approximated
from stored teacher hidden states. Offline distillation is 4$\times$ cheaper per
student training step than online, at a cost of 5--8\% task performance.

\subsection{Online Distillation}

In online distillation, the teacher model runs synchronously during student training,
providing fresh soft labels and hidden states for each training batch. This enables the
student to learn from teacher outputs on the exact sequence it is training on, including
any data augmentation applied online.

\subsection{Progressive Distillation}

Progressive distillation \cite{hinton2015distilling,salimans2022progressive} applies a
chain: $32\text{B} \to 8\text{B} \to 4\text{B} \to 0.6\text{B}$, where each stage
uses the previous stage's output as the teacher. This is more effective than direct
compression from 32B to 0.6B because each step makes a manageable capacity reduction.

\subsection{Task-Specific Distillation}

For deployment scenarios where only a specific task matters (e.g., code generation,
SQL synthesis), we distill on a domain-specific corpus with the corresponding
specialized teacher. Task-specific students achieve near-specialist performance in
their domain while fitting in an 8B parameter budget.

%% -----------------------------------------------------------------------
\section{Experiments}
\label{sec:experiments}
%% -----------------------------------------------------------------------

\subsection{Teacher-Student Performance Gaps}

\begin{table}[H]
\centering
\caption{Performance gap closure (\%) for different distillation methods.
Teacher = Zen-32B (Qwen3-32B), student = Zen-8B (Qwen3-8B).
Gap closure = (student - base) / (teacher - base) $\times$ 100, i.e.\ the fraction of the
8B$\to$32B benchmark gap recovered by distillation. Values are relative to the undistilled
8B base and compare distillation \emph{methods}; they are not standalone capability claims.}
\begin{tabular}{lrrrr}
\toprule
\textbf{Method} & \textbf{MMLU} & \textbf{GSM8K} & \textbf{HumanEval} & \textbf{Avg.} \\
\midrule
No distillation (8B base) & 0\%  & 0\%  & 0\%  & 0\% \\
Standard KD               & 58\% & 61\% & 64\% & 61\% \\
SPD (ours)                & 72\% & 76\% & 73\% & 74\% \\
SPD + multi-teacher       & 74\% & 81\% & 79\% & 78\% \\
\bottomrule
\end{tabular}
\label{tab:gap_closure}
\end{table}

\subsection{Effect of Distillation Method on Student Quality}

Holding the student fixed at Zen-8B (Qwen3-8B) and varying the distillation method, we
observe a consistent ordering on MMLU, GSM8K, and HumanEval: the undistilled base is
weakest, standard KD improves it, SPD improves further, and SPD with multi-teacher
aggregation is strongest, approaching the Zen-32B teacher. The same method ordering holds
for the progressively-distilled 0.6B student, which recovers a smaller but still
substantial fraction of teacher quality at a much smaller parameter budget.

\subsection{Distillation Efficiency}

The motivation for distillation is the inference-cost/quality trade-off: the smaller
distilled students run several times faster than the 32B teacher while retaining most of
its benchmark quality. The Zen-8B SPD+MT student offers a favorable operating point
(near-teacher MMLU at a multiple of the teacher's throughput), and the progressively
distilled 0.6B student offers the highest throughput for the most cost-sensitive
deployments, at a larger quality reduction. Exact throughput depends on hardware and
serving configuration.

\subsection{Scaling Laws for Distillation}

We observe that distillation efficiency (performance per parameter) follows a power law
in the student-to-teacher size ratio:

\begin{equation}
    \text{Gap closure}(s) = 1 - C \cdot \left(\frac{s}{t}\right)^{-\delta}
    \label{eq:scaling}
\end{equation}

where $s$ is the student parameter count and $t$ is the teacher parameter count. Fitting
this form across the Qwen3-based Zen student sizes (0.6B, 4B, 8B) distilled from the
Zen-32B teacher, gap closure increases smoothly and concavely with $s/t$: larger students
recover a larger fraction of the teacher's quality, with diminishing marginal returns as
the student approaches the teacher's size. The fitted curve lets practitioners estimate
the gap closure achievable at a target student size before committing to a distillation
run.

\subsection{Layer Matching Ablation}

Ablating the SPD layer-matching strategy on the Zen-8B student, adding any representation
alignment improves over KD alone, and the choice of \emph{which} layers to align matters:
probing-based matching (aligning the few teacher layers most predictive of task accuracy)
performs best, uniform matching (every few layers) is close behind, and aligning all
layers performs slightly worse than the targeted strategies — indicating that a small
number of well-chosen alignment points captures most of the benefit.

%% -----------------------------------------------------------------------
\section{Task-Specific Distillation Results}
\label{sec:task_specific}
%% -----------------------------------------------------------------------

Comparing task-specific Zen-8B (Qwen3-8B) students against the general SPD+MT Zen-8B
student on four domain benchmarks (HumanEval, GSM8K, LegalBench, MedQA), each
domain-specialized student is the strongest on its own target benchmark: the code-specific
student leads on HumanEval, the math-specific student on GSM8K, the legal-specific student
on LegalBench, and the medical-specific student on MedQA. Specialization trades some
general breadth for a clear gain in the target domain, recovering close to specialist-level
teacher performance within an 8B parameter budget.

%% -----------------------------------------------------------------------
\section{Conclusion}
\label{sec:conclusion}
%% -----------------------------------------------------------------------

Semantic-preserving distillation (SPD) with multi-teacher aggregation closes a
substantially larger fraction of the 32B-to-8B performance gap than standard KD (78\% vs.\
61\% of the gap in our experiments). Progressive distillation enables a large throughput
increase at the 0.6B scale. The distillation scaling law (Equation~\ref{eq:scaling})
predicts gap closure as a function of $s/t$, enabling practitioners to select the optimal
student size for their compute budget. Task-specific distillation recovers close to
specialist-level teacher performance in target domains, providing an efficient deployment
path for specialized applications. All teachers and students are Qwen3-based Zen models
(Apache-2.0).

\begin{thebibliography}{9}
\bibitem{hinton2015distilling}
G. Hinton, O. Vinyals, J. Dean.
\textit{Distilling the Knowledge in a Neural Network}.
arXiv:1503.02531, 2015.

\bibitem{salimans2022progressive}
T. Salimans, J. Ho.
\textit{Progressive Distillation for Fast Sampling of Diffusion Models}.
ICLR, 2022.

\bibitem{park2019relational}
W. Park, D. Kim, Y. Lu, M. Cho.
\textit{Relational Knowledge Distillation}.
CVPR, 2019.

\bibitem{romero2015fitnets}
A. Romero, N. Ballas, S.E. Kahou, et al.
\textit{FitNets: Hints for Thin Deep Nets}.
ICLR, 2015.
\end{thebibliography}

\end{document}
